\documentclass[fontset=none,aspectratio=1610]{ctexbeamer}

\usepackage{minted}
\usepackage{graphicx}
\usepackage{fontspec}
\usepackage{amssymb, amsmath, amsfonts}

\usefonttheme{serif}
\usetheme{CambridgeUS}

\setmainfont{Noto Serif CJK SC}
\setsansfont{Noto Sans CJK SC}[Scale=MatchLowercase]
\setmonofont{Noto Sans Mono CJK SC}[Scale=MatchLowercase]
\setCJKmainfont[ItalicFont = FZKai-Z03]{Noto Serif CJK SC}
\setCJKsansfont[ItalicFont = *]{Noto Sans CJK SC}
\setCJKmonofont[ItalicFont = *]{Noto Sans Mono CJK SC}

\setlength{\parskip}{0.5em}

\setminted{
    baselinestretch=1.0,       % 行间距
    frame=single,           % 边框样式
    framesep=2mm,              % 边框间距
    rulecolor=\color{black},   % 边框颜色
    breaklines=true,           % 自动换行
    breakanywhere=false,       % 不在任意位置断行
    showspaces=false,          % 不显示空格
    showtabs=false,            % 不显示制表符
    tabsize=4,                 % 制表符大小
}

\title{第六次上机课}
\author{臧炫懿}
\date{\today}

\begin{document}
\frame{\titlepage}
\begin{frame}
    \frametitle{目录}
    \tableofcontents
\end{frame}

\begin{frame}
    \frametitle{本节课同学们的目标}

    听故事（划去）

    这节课的内容和课程本身关系不大，这完全可以归因于这周没有正课。而且我也不知道期中考试是哪一天，如果我知道的话就带着大家复习了。

    本节课的内容可以简单理解为你北计算机神课ICS+四大礼包《操统》《编原》《计网》《体系》的超级蒸馏简化版，虽然内容不多，但我觉得对你们理解计算机的整体架构和运作机制会有很大帮助。

\end{frame}

\section{现代计算机体系结构}
\frame{\sectionpage}
\subsection{硬件浅论}
\begin{frame}
    \frametitle{现代计算机的硬件组成}

    \begin{description}
        \item[计算设备] 包括CPU这个主要的通用计算设备，以及GPU等其他的专门计算设备，也包括PCH芯片组（南桥）等一些辅助计算设备。
        \item[内存] 包括RAM和ROM等。
        \item[外存] 包括硬盘、SSD等。
        \item[主板] 是连接各个硬件组件的核心电路板，所有的硬件组件都通过主板进行连接和通信。这些通信方式包括总线、接口和协议等。
        \item[输入输出设备] 包括键盘、鼠标、显示器等。
        \item[电源] 电源不参与计算或数据存储，但它为计算机的各个组件提供必要的电力支持，使得计算机能够正常运行。
    \end{description}

\end{frame}

\begin{frame}
    \frametitle{现代计算机硬件体系结构}
    
    现代计算机硬件的体系结构大致是这样组织的：
    \begin{itemize}
        \item \textbf{CPU}是计算机的核心处理器，负责执行指令和处理数据。
        \item CPU通过\textbf{内存控制器}和\textbf{内存总线}直接访问并控制\textbf{主存}。
        \item CPU通过\textbf{PCIe控制器}和\textbf{PCIe总线}连接到其他的高速设备，如\textbf{GPU}、\textbf{PCH芯片组（南桥）}、\textbf{NVMe协议的SSD}等。
        \item 南桥芯片通过被统称为\textbf{I/O总线}的接口连接到各种输入输出设备和其他低速设备，如\textbf{键盘}、\textbf{鼠标}、\textbf{显示器}、\textbf{USB设备}、\textbf{SATA协议的SSD和HDD}等。
        \item 现代高性能显卡也可以直接连接到显示器上，而不需要经过南桥芯片和I/O总线。
    \end{itemize}

\end{frame}

\subsection{计算机怎么启动}

\begin{frame}
    \frametitle{按下电源键之前}

    研究计算机的启动过程，可以帮助我们更进一步地理解计算机的体系结构和运作机制。

    计算机在关机状态下，绝大多数硬件确实是完全断电状态。但反直觉的是，这段时间内计算机仍是通电状态（由CMOS电池提供电力）。
    
    CMOS供电的一个重要部件是BIOS/UEFI芯片，记录了当前时间、日期、硬件配置等信息。如果抠掉该电池，计算机虽然大概率仍能启动，但时间等信息会丢失。

\end{frame}

\begin{frame}
    \frametitle{按下电源键之后}

    \begin{enumerate}
        \item 按下电源键时，主板收到信号，并给电源发信号。
        \item 电源率先启动，开始按一定的时序输出指定电压并自检。自检完毕，发出Power Good信号，通知主板电压稳定。
        \item 主板收到信号，芯片组开始工作，复位CPU。
        \item CPU复位后，进入\textbf{实模式}，开始运行第一个程序：BIOS/UEFI程序。
        \item POST。
        \item 其余硬件次序启动，为PCIe设备分配总线编号、内存映射IO空间、中断号等。
    \end{enumerate}

    自此所有硬件准备完毕，开始试图进入引导阶段。

\end{frame}

\begin{frame}
    \frametitle{引导阶段}

    \begin{enumerate}
        \item BIOS/UEFI程序开始按预定顺序寻找引导设备。旧的BIOS程序会寻找MBR，现代的UEFI程序会寻找EFI系统分区。没找到则发出错误信号，并终止启动流程。
        \item 如果找到了操作系统，B/U程序会把控制权交给操作系统的引导程序。
        \item 引导程序继续寻找操作系统内核，并把控制权和必要的参数交给内核。
        \item 内核获取控制权之后，初始化CPU调度子系统、内存管理子系统、硬件抽象层、驱动程序。
        \item 下一步，第一个用户态程序被加载并执行（ \texttt{smss.exe} 、 \texttt{init} 、 \texttt{systemd} 等）。用户态程序继续加载其他的用户态程序，直到系统完全启动。
    \end{enumerate}

\end{frame}

\subsection{计算机怎么输入输出}

\begin{frame}
    \frametitle{输入输出设备的工作原理}

    上述设备其实都是系统设备，一个不能与外界交互的裸机也是这么启动的。但不能交互的裸机用处较为有限，我们必须有输入输出设备才能最大发挥计算机的功能。

    MMIO是现代计算机最常见的输入输出方式。CPU不认识谁是谁，它只把这些地址当成普通内存读写，然后由主板芯片组去翻译成PCIe事务并送到目标设备。所以同一条 \texttt{mov} 指令，可能是访问内存，也可能是访问设备，这取决于地址。

\end{frame}

\begin{frame}
    \frametitle{OS到底干了一个什么事}

    在进行以下的讲解之前，我们先要弄明白这个问题。实际上，很难对OS下一个准确的定义，因为OS的功能实在是太多了，甚至有些功能是不同的OS之间差异很大的。但我们可以从一个比较抽象的角度来理解OS：OS是一个管理计算机资源和提供服务的软件系统。

    实际上OS的功能是极多的，但我们总体上把这堆功能分到两大类：
    \begin{description}
        \item[资源管理] 包括CPU调度、内存管理、文件系统、设备管理等。这是所有OS都必须具备的核心功能。
        \item[服务提供] 包括系统调用、网络服务、安全服务、命令行接口、图形用户界面等。这些功能可能在不同的OS之间差异很大，但它们都是为了让用户能够更方便地使用计算机而提供的服务。
    \end{description}

    以下内容就围绕着这两个方面来展开。

\end{frame}

\begin{frame}
    \frametitle{轮询、中断和DMA}

    最朴素的输入输出方式是轮询，CPU不断地检查设备状态寄存器，看看设备是否准备好进行数据传输。但这样，效率会极为低下。

    为解决这个问题，现代计算机引入了中断机制。CPU会做自己的事情，当设备准备好进行数据传输时，会沿着中断请求线发出一个信号，中断控制器收到信号后通知CPU，CPU暂停当前工作，转而执行中断处理例程来处理这个事件。中断机制大大提高了输入输出的效率。

    但中断只能通知“有事”，不能直接进行数据传输。DMA就派上了用场，它常用于现代SSD或网卡等数据量较大的设备。这时候，CPU只需要告诉DMA控制器数据的来源、目的地和大小，DMA控制器就可以直接在内存和设备之间进行数据传输，而不需要CPU的干预。

\end{frame}

\begin{frame}
    \frametitle{驱动}

    我们反复提及“驱动”。

    刚刚说了，有了MMIO、轮询中断DMA这堆东西，CPU就能够和外设打交道了。但每个外设的工作方式是不一样的，CPU不可能为每个外设都写一套代码来和它打交道。
    
    驱动程序就是为每个外设编写的一套软件接口，抽象了设备的细节，让操作系统和应用程序能够以统一的方式与各种不同的设备进行交互，即“同声传译”。

    在启动流程中，我们提及驱动还要管设备的电源，一个驱动bug就会导致设备“睡死过去”，所以一般认为“驱动占了70\%的代码和90\%的bug”。

\end{frame}

\begin{frame}
    \frametitle{热插拔和总线枚举}

    早年的打印机等必须关机才敢插拔（冷插拔），因为插拔的时候容易产生电压波动，可能会损坏设备或主板。现代计算机支持热插拔，允许用户在计算机运行时插拔设备而不会造成损坏。这得益于接口处加上的电气隔离电路，以及操作系统的总线枚举机制。

    所谓总线枚举，可以理解为：
    \begin{enumerate}
        \item 新设备插入，接口芯片检测到电压变化，发出信号“present”。
        \item 操作系统收到信号，开始在总线上广播“你是什么东西？”，设备回复自己的ID、类型、功能等信息。
        \item 操作系统根据设备信息，加载相应的驱动程序，并完成设备的初始化和配置，以及资源的分配。
    \end{enumerate}

\end{frame}

\begin{frame}
    \frametitle{实例}

\begin{enumerate}
\item 用户点“保存”，应用调用 \texttt{write()} 系统调用，内核把用户缓冲区映射到页缓存；
\item 文件系统试图把页缓存写回磁盘，发现是个U盘，调用U盘驱动做“写”这个事；
\item 驱动把 SCSI 命令封装成 CBW，交给主机；
\item 主机把 CBW 切成微帧，通过差分信号线发送；
\item U 盘主控收到 CBW，把逻辑块地址翻译成 NAND 物理页；
\item 主机发起 DMA，把数据突发写入 U 盘内部缓存，U 盘回传“写完”；
\item 内核收到“写完”，把页缓存标记为干净，应用弹出“保存成功”。
\end{enumerate}

\end{frame}

\section{计算机操作系统}

\subsection{信息如何被表示}

\begin{frame}
    \frametitle{二进制和编码}

    计算机只看得懂二进制，所以所有的信息都必须被编码成二进制的形式。对于不同类型的信息，计算机使用不同的编码方式来表示它们。

    \begin{description}
        \item[定点数] 二进制补码。
        \item[浮点数] IEEE 754。
        \item[字符] ASCII、Unicode等。 
    \end{description}

    \textbf{必须说明，计算机本身并不知道一段特定的二进制数据代表什么信息，这完全取决于我们如何解释它。} 于是就知道了为什么C语言中的 \texttt{union} 可以让同一段内存既可以被解释为整数，也可以被解释为浮点数了。

\end{frame}

\begin{frame}
    \frametitle{二进制补码}

    可以理解为：除了最高位是负的，其余都是正的。

    比如 \texttt{INT4} （在模型训练中常用），从右往左分别代表：$2^0$、$2^1$、$2^2$、$-2^3$。显然 \texttt{INT4} 的范围是$[-8, 7]$。
    
    这带来相当的好处：
    \begin{itemize}
        \item 取符号：直接看最高位，用掩码表示就是 \texttt{AND 1000}。
        \item 取相反数：直接取反再加1，用掩码表示就是 \texttt{NOT} 再 \texttt{ADD 1}。证明略，应当不难。
        \item 取绝对值：如果是负数，先取反再加1；如果是正数，保持不变。
        \item 加法和减法可以用同一套电路实现，减法可以看成加上一个负数。
        \item 避免了$+0$和$-0$的存在。
    \end{itemize}

\end{frame}

\begin{frame}
    \frametitle{半加器}

    半加器的真值表：
    \begin{table}[h]
        \centering
        \begin{tabular}{cc|cc}
            \hline
            A & B & S & C \\
            \hline
            0 & 0 & 0 & 0 \\
            0 & 1 & 1 & 0 \\
            1 & 0 & 1 & 0 \\
            1 & 1 & 0 & 1 \\
            \hline
        \end{tabular}
    \end{table}
    其中S是A和B的异或（XOR），C是A和B的与（AND）。全加器在半加器的基础上增加了一个进位输入Cin，输出S和进位输出Cout。学过《微电基》的同学应该对它们不陌生了。
\end{frame}

\begin{frame}
    \frametitle{IEEE 754}

    IEEE 754是浮点数的国际标准，定义了浮点数的表示方法和运算规则。一个浮点数包含三个部分：符号位（S）、指数位（E）和尾数位（M）。以单精度浮点数为例，S占1位，E占8位，M占23位。该浮点数的值可以表示为：
\[(-1)^S \times 1.M \times 2^{E - 127}\]

    其余的浮点数类型（如双精度、半精度等）也是类似的，只是S、E、M的位数不同。

    现在的问题就变成：尾数位如果全是0，指数位如果全是0或全是1，这些特殊情况该怎么处理？IEEE 754做了特别规定：
\begin{itemize}
    \item 当E全为0且M全为0时，表示0。
    \item 当E全为0且M不全为0时，表示非规格化数。
    \item 当E全为1且M全为0时，表示无穷大。
    \item 当E全为1且M不全为0时，表示NaN（Not a Number）。
\end{itemize}

\end{frame}

\begin{frame}
    \frametitle{IEEE 754中的非规格化数}

    非规格化数是指当指数位全为0时，尾数位不全为0的情况。在单精度浮点数中，非规格化数的值可以表示为：
\[(-1)^S \times 0.M \times 2^{-126}\]
    其中M是尾数位的二进制小数部分。非规格化数的引入是为了表示非常接近于0的数值，避免了在规格化数中出现的下溢问题。

    于是我们发现了一个问题：\textbf{浮点数的本质是科学记数法}。

\end{frame}

\begin{frame}
    \frametitle{浮点数的警示}

\begin{itemize}
    \item 浮点数的精度是有限的，不能表示所有的实数。
    \item 浮点数不是稠密的。
    \item 浮点数的运算是先将两个数转换成相同的指数，然后进行尾数的加减乘除，最后再进行规格化处理。这可能会导致精度损失和舍入误差，所以不满足加法和乘法的交换律和结合律。
\end{itemize}

因此\textbf{不得用浮点数进行货币计算}（除非四精度 \texttt{decimal} 或专门的库如 \texttt{libmpdec} ）。常规的表示实际是定点数分币法。

另外，在浮点数运算的时候，尽量保证大小相近的数先运算，再进行数量级差异较大的数的运算，以减少精度损失。

\end{frame}

\begin{frame}
    \frametitle{字符：编码方案的漫漫统一路}

    我们知道，现在的计算机是使用二进制来存储数据的。但是，对于人类使用的语言而言，即使是相对简单的英语，也有52个大小写字母、10个数字、各种标点符号和特殊符号等，远远超过了二进制的0和1两种状态。

    容易想到的做法是：给每一个字符分配唯一的一个编号，然后再使用二进制来表示该编号。这样的一个单射就叫做“编码方案”。

    但一开始不同厂家用的是不同的编码方案，导致同一个编码在不同的计算机上表示出不同的字符，对数据交换造成了极大的麻烦，如“暗黑魔法师”变成了“穞堵臸猭畍”。

\end{frame}

\begin{frame}
    \frametitle{要么解决问题，要么解决提出问题的人}

    为了解决上述问题，显然有两个方案：要么利用一些方法区分不同的编码方案，并在使用时现场翻译（解决问题），要么强制所有人都使用同一个编码方案（解决提出问题的人）。但前者太麻烦了，所以还不如解决提出问题的人（谁家秦始皇统一度量衡，谁家乔布斯统一接口）。

    于是美国人率先制定了ASCII编码方案，使用7位二进制来表示128个字符，满足了英语的需求。但对于其他语言来说，128个字符远远不够用，例如法语、德语等欧洲语言需要更多的字符来表示各种变音符号和特殊字符。于是又出现了ISO 8859-1（Latin-1）、扩展ASCII等编码方案，使用8位二进制来表示256个字符，满足了更多语言的需求。

\end{frame}

\begin{frame}
    \frametitle{汉语：你牛，但穞更牛}

    但是当涉及到汉语时，情况就变得非常复杂了。汉字的数量远远超过了256个，甚至翻一百倍也不够用。

    我国也做过相关的努力，如GB2312，后来又发展出了GBK等。但依然需要考虑诸如日本、韩国、越南等国家、中国台湾等地区的不同需求，这些国家和地区的当局也开发了自己的编码方案，如Shift-JIS、EUC-KR、Big5等。

    结论：感觉不如世界人民大团结。

    于是，Unicode诞生了。

\end{frame}

\begin{frame}
    \frametitle{Unicode}

    Unicode使用16位二进制来表示65536个字符，满足了全球各种语言的需求，甚至还有富余，能容纳更多的符号甚至表情符号（emoji）等。但是实际上65536也远远不够：如果再向上追溯，光是中国古代弃用的那些“死汉字”也有几万个，甚至几十万个。于是Unicode又引入了代理对（surrogate pair）的机制，使用两个16位的编码单元来表示一个字符，从而扩展了Unicode的编码范围。

    Unicode有不同的表示方式，如UTF-8、UTF-16、UTF-32等。Linux和macOS等操作系统通常使用UTF-8编码（不兼容GBK），而Windows则使用UTF-16编码（但中国地区的计算机依然是GBK和UTF-16并存），于是又出现了新的问题，真实一头烂帐。好在Windows一直在测试UTF-8，未来可能会统一到UTF-8上。

\end{frame}

\begin{frame}
    \frametitle{从字符到字形和字体}

    Unicode算是为每一个抽象字符分配了一个唯一的编码，如汉字“汉”的Unicode编码是U+6C49，但这仅仅是规定了U+6C49这个编码代表了“汉”这个抽象字符，并没有规定这个字符应该如何显示出来。

    真正的字符长相，由\textbf{字形}决定，如宋体、黑体、楷体、仿宋都是汉字的不同字形。把目光放远到全世界，我们发现同一个字符在不同的国家和地区也有明显差异，例如中国大陆的汉字和中国香港、中国台湾、日本、韩国等地的字形都有明显的差异。

    将一套风格相同的字形放在一起，装进一个索引文件，就形成了一个\textbf{字体文件}。一组字体文件就构成了一个\textbf{字体族}，简称“字体”。如宋体、黑体、楷体、仿宋等都是常见的中文字体族。

\end{frame}

\begin{frame}
    \frametitle{从字体到屏幕}

    计算机屏幕是无数个像素点组成的，所以要把一个字符显示在屏幕上，就需要把这个字符的字形转换成像素点的形式，即“字体渲染”。

    朴素的做法就是把字体中的每一个字形作为图片来存储，即所谓的“位图”字体。但这样就会出现一个弊病：如果我们想要一个大号字，简单放大图片就会变得模糊不清；如果我们想要一个小号字，简单缩小图片就会变得难以辨认。总不能为每一个字符都存储多个大小的图片吧，这样会占用大量的存储空间。

\end{frame}

\begin{frame}
    \frametitle{王选的点阵字模，PS、TTF和OTF}

    1978年，王选院士发明了点阵字模技术，使用一个矩阵来表示一个字符的字形。这样可以轻易地生成不同字号的汉字，还数千倍地压缩了存储空间，突破了位图字体的局限性，使字符能被“算出来”而不是“画出来”，彻底打开了人们编写字体的思路。

    1982年，Adobe推出划时代的字体文件PostScript，用三次贝塞尔曲线描述字形，被称作“矢量字体”，比点阵字模技术更进一步。后来苹果搞出了TrueType字体，微软和Adobe又联手搞出了OpenType字体，这些都是基于矢量字体的技术，支持更丰富的字形特征和更高的显示质量。

    时至今日，汉字也依然是大字号矢量字形、小字号点阵字模的混合使用，来兼顾显示质量和存储空间，王选院士也被誉为当代毕昇。

\end{frame}

\begin{frame}
    \frametitle{字重、字宽、斜体等特征}

    在计算机出现之前，这些名词就随着印刷术的发展而出现了。既然把字体搬到计算机上，自然会考虑设计不同的字重和斜体版本。这些都是字体的特征，也可以不设计，但设计了就可以让用户有更多的选择。

    然而，大多数汉字字体并没有斜体字体：这是因为汉字是方块字，斜着不好看，没有这方面的需求，因此大多数汉字字体并没有斜体版本，通常用楷体来表示斜体的效果。对于英文字体而言，斜体是非常常见的，例如Times New Roman、Arial等都有斜体版本。

\end{frame}

\begin{frame}
    \frametitle{衬线、无衬线和等宽}

    字体可以分为衬线、无衬线两大类。

    衬线字体是指在字形的笔画末端有装饰性的小线条，这些小线条被称为“衬线”。这种字体用于印刷和长篇阅读，例如Times New Roman、Georgia、宋体等。

    \textsf{无衬线字体在字形的笔画末端没有衬线，更多用于屏幕显示和标题，例如Arial、Helvetica、黑体等。}

    有一种特殊的字体被称为“等宽”字体，也叫打字机字体，指的是每个字符占用相同宽度的字体。这种字体常用于编程和终端显示，例如Courier New、Consolas、仿宋等。

    汉字天然是方块字，所以天然等宽，所以为汉字设计等宽字体没什么意义。通常利用仿宋（搭配衬线等宽）或黑体（搭配无衬线等宽）来达到等宽的效果。即使是为汉字设计等宽字体，也实际上是为字体族中的西文字符设计了相当于半个汉字宽度的等宽版本。

\end{frame}

\begin{frame}
    \frametitle{LF和CRLF}

    \texttt{CR} 是“回车”，控制字符是 \texttt{\textbackslash r}；\texttt{LF} 是“换行”，控制字符是 \texttt{\textbackslash n}。现代计算机上这两个是一回事，但打字机时代回车是把打字头回到行首，换行是把纸往下送一行，所以它们是两个不同的动作。那为什么现代计算机上只剩下 \texttt{Enter} （回车）了呢？归根结底其实是因为“麻烦”。

    在上世纪，Unix用 \texttt{LF}，Dos用 \texttt{CRLF}，Mac用 \texttt{CR}。后来Mac内核换成了BSD（mac OS X），也就改成了 \texttt{LF}；Linux直接继承Unix，也就用 \texttt{LF}；Windows则继承了Dos（因为Windows的成名作Win3.1就是Dos上的一个软件），于是就继续用 \texttt{CRLF} 了，即使现在的Windows和Dos已经没什么关系了也没有改过来。

\end{frame}

\subsection{程序怎么跑起来}

\begin{frame}
    \frametitle{从代码到二进制}

    我们写的代码，充其量只是一堆有意义的文本。计算机只认二进制，不认文本，那代码是怎么跑起来的呢？

    这需要一个东西充当翻译的功能。对于C/C++，需要从源代码经过\textbf{编译}变成目标文件，再经过\textbf{链接}变成可执行文件；对于Python，需要从源代码经过\textbf{编译}变成字节码，再由Python虚拟机解释执行。

    前者叫编译型语言，后者叫解释型语言。

\end{frame}

\begin{frame}
    \frametitle{CPU计算原理}

    CPU内部，去掉后来添加的一堆花里胡哨的东西（如被吃掉的北桥、内存控制器），主要由算术逻辑单元（ALU）、寄存器、控制单元等组成。CPU通过执行指令来进行计算，每条指令都包含一个操作码（opcode）和一些操作数（operands）。操作码指定了要执行的操作类型，而操作数则提供了执行该操作所需的数据。

    寄存器是CPU内部的高速存储器，用于存储指令执行过程中需要使用的数据和地址。ALU负责执行算术和逻辑运算，如加法、减法、乘法、除法、与、或、非等。控制单元则负责协调CPU内部各个组件的工作，确保指令按照正确的顺序执行，并处理各种异常情况。

    人们把指令写成这样的形式： \texttt{mov \%rax \%rbx} ，这样的指令叫做“汇编”。汇编和机器码只差一步，所以很多语言都是先编译成汇编，再由汇编器把汇编翻译成机器码的。这里的 \texttt{mov} 是操作码， \texttt{\%rax} 和 \texttt{\%rbx} 是操作数。CPU执行这条指令时，会把寄存器 \texttt{rax} 的值复制到寄存器 \texttt{rbx} 中。

\end{frame}

\begin{frame}
    \frametitle{C语言编译流程}

    \begin{description}
        \item[预处理] 对源代码进行文本插入（ \texttt{\#include} ）、宏展开和宏替换（ \texttt{\#define} ）、条件编译（ \texttt{\#ifdef} ）、移除注释等处理，生成预处理文件。这步得到的代码和源代码其实是完全等价的。
        \item[编译] 编译器把预处理文件翻译成汇编代码，将源代码中的每个函数、变量等转换为对应的汇编指令和数据定义，并生成符号表来记录这些符号的信息。同时也会进行编译时优化。\textbf{优化会暴露UB}。
        \item[汇编] 汇编器把汇编代码翻译成机器码，生成目标文件，包含了机器码、符号表、重定位信息等。
        \item[链接] 链接器，把所有的目标文件和库文件中的符号表合并，进行符号解析和重定位，最终生成一个包含所有代码和数据的可执行文件。 
    \end{description}

\end{frame}

\begin{frame}
    \frametitle{其他编译情况}

    \begin{description}
        \item[Python] 解释性语言，先由解释器（CPython）把源代码编译成字节码（.pyc文件），再由虚拟机（Python VM）解释执行字节码。另一些工具（PyPy、Jython）则使用即时编译（JIT）技术，在运行时将字节码编译成机器码以提高性能。
        \item[C\#] “中间语言”。先由编译器（Roslyn）把源代码编译成中间语言（IL），再由公共语言运行时（CLR）使用即时编译（JIT）技术将IL编译成机器码并执行。分发时，大多直接分发IL代码，所以C\#的逆向非常简单。
    \end{description}

    不谈C++、Rust等。拿这两个玩意讲编原，那是嫌自己活得不耐烦了。前者的异常处理、虚函数表、模板，后者的所有权系统、trait对象等，都是编译器的重口味大餐，吃不吃得下就看个人了。我虽然会，但我不能在这种课上说。

\end{frame}

\begin{frame}
    \frametitle{从可执行文件到程序}

    在Linux中，可执行文件叫做ELF；在Windows中，可执行文件叫做PE。

    运行一个可执行文件时，OS通过加载器把可执行文件加载到内存中，并进行必要的重定位和符号解析，创建一个新的\textbf{进程}来运行这个程序。进程中可能有多个\textbf{线程}，每个线程都有自己的执行上下文（寄存器、堆栈等），但它们共享同一个进程的内存空间。\textbf{进程是资源分配的基本单位，线程是CPU调度的基本单位}。

    程序在执行的过程中，CPU会不断地从内存取指令和数据，并次第执行这些指令。执行完毕后，OS会回收进程占用的资源，结束程序的运行。程序本身大多无法直接操作分配给自己内存以外的资源，必须通过系统调用来请求OS提供服务，如文件操作、网络通信、进程管理等。

\end{frame}

\begin{frame}
    \frametitle{stack和heap}

    程序运行时，内存通常被划分为几个不同的区域，其中最重要的两个区域是\textbf{栈（stack）}和\textbf{堆（heap）}。前者由编译器自动管理，小而快；后者由程序员手动管理，大而慢。

    所有的 \texttt{malloc} 都是在堆上分配内存的。而诸如 \texttt{int[100] a} 这类局部变量内存分配是在栈上分配的。特殊的是 \texttt{static int a} 这类静态变量内存分配是在数据段分配的，数据段和堆栈是完全不同的两个区域。

    栈虽然快，但资源往往很紧张，尤其是在使用递归调用的情况下，很容易导致堆栈溢出，俗称“爆栈”。从安全性上考量，这给了我们两个启示：
    \begin{itemize}
        \item 多 \texttt{malloc} ，少 \texttt{int[100]} 。
        \item 递归调用最好换个写法，例如迭代等。
    \end{itemize}

\end{frame}

\begin{frame}
    \frametitle{通过反汇编学习程序运行机制}

    用C随便写点东西，然后用 \texttt{gcc} 编译（不要用 \texttt{g++} ），再用 \texttt{objdump -d} 反汇编，就可以看到程序是怎么运行的了。

    试着找一找程序的代码段（ \texttt{.text} ）、数据段（ \texttt{.data} ）、BSS段（ \texttt{.bss} ）。

\end{frame}

\begin{frame}
    \frametitle{链接原理}

    在编译的时候，对每一个目标文件，都会生成一个符号表，记录了这个目标文件中定义的符号（函数、变量等）和引用的符号（外部函数、变量等）。在这一段，不会出现任何有关“符号未定义”的报错：编译器默认这堆东西在别的目标文件里定义了，先放一边，等链接的时候再说。而这些符号也都会记录其地址，这些地址是基于当前目标文件的偏移量。

    而链接实际上就干了一个事：把这堆符号表合并成一个全局符号表。在链接的时候，链接器会：
    \begin{itemize}
        \item 把这些目标文件的代码段和数据段合并成一个连续的内存区域。
        \item 对这些目标文件的符号表进行处理：代码段中的相对偏移量被替换为基于全局符号表的全局偏移量，即 \textbf{重定位} 。
        \item 如果出现重声明（一般不允许）、重定义（新标准不允许）或未定义（不允许）的符号，链接器会报错。
    \end{itemize}

\end{frame}

\begin{frame}
    \frametitle{特殊变量的符号表}

    我们知道，C语言的变量按作用域和生命周期等可以分为不同的类型。

    \begin{description}
        \item[普通的全局变量] 在符号表中会有一个全局符号，指向数据段中的一个地址。
        \item[普通的局部变量] 局部变量在符号表中不会有任何符号，因为它们是在栈上分配的，地址是动态的，编译器会通过寄存器或栈指针来访问它们。
        \item[static的全局变量] 静态的全局变量，在全局符号表中会有一个符号，但这个符号是局部的（local），只能在当前目标文件中访问，链接器不会去参与全局符号表的解析。这是非常难以理解的一个特例。
        \item[static的局部变量] 局部变量，就别想在全局符号表里找到它了。
        \item[extern的变量] 意思是“这个变量在别的目标文件里定义了”，所以自然会在全局符号表中出现一个符号，但这个符号是未定义的（undefined），链接器会在链接的时候去寻找这个符号的定义。 
    \end{description}
    
\end{frame}

\subsection{内存怎么被管理}

\begin{frame}
    \frametitle{那么计算机是怎么管理上面这堆机器码的？}

    我们知道，计算机的CPU从内存取取指令和数据来执行程序，并把结果写回内存。

    但是对于一些用户，我们可能会在后台挂着114514个进程，这些进程都要吃内存。这些进程需要的总内存可能远远大于实际物理内存的大小，那么计算机是怎么管理这些内存，使得它们能够正常运行的呢？

\end{frame}

\begin{frame}
    \frametitle{虚拟内存}

    计算机使用虚拟内存来管理内存，每个进程都有一个独立的虚拟地址空间。所以对于任何程序而言，他们的指针都是从0x00000000开始的，完全不需要关心其他程序的内存地址！而OS会把这些虚拟地址映射到实际的物理内存地址上，这个过程叫做地址转换。

    对于普通程序开发者，我们完全不需要关心物理内存的地址，只需要使用虚拟地址就可以了。

    而这时，我们会发现虚拟内存还能达成一个意外之喜：虚拟内存的存在，实际上阻止了程序访问不属于它的内存区域。因为每个进程都有自己的虚拟地址空间，操作系统只需确保每个进程只能访问自己的虚拟地址空间，不能访问其他进程的虚拟地址空间。这样就能有效地保护每个进程的内存安全，防止恶意程序破坏系统的稳定性和安全性。这是现代操作系统比DOS等老式操作系统更安全的重要原因之一。

\end{frame}

\begin{frame}
    \frametitle{页}

    操作系统管理内存的基本单位是页（page）。虚拟内存和物理内存都被划分成固定大小的页，OS实际上就是通过页面的映射来实现虚拟内存和物理内存之间的地址转换的。

    具体的实现方式是通过一个叫“页表”的数据结构来记录虚拟地址和物理地址之间的映射关系。当CPU访问一个虚拟地址时，MMU会查找页表，找到对应的物理地址，然后进行访问。

\end{frame}

\begin{frame}
    \frametitle{swap和缺页异常}

    而真正解决物理内存不够问题的是磁盘交换区（swap）。当物理内存不足时，OS会把一些不常用的页从DRAM移到swap上，这样就腾出物理内存来给其他进程使用。

    这就会导致一个问题：如果CPU想要访问一个虚拟地址，MMU去查页表，发现这个页不在DRAM中，这时候就会发生一个缺页异常。OS就会捕获这个异常并去swap上找到这个页，把它加载回DRAM中，然后更新页表，最后让CPU重新访问这个虚拟地址，这样就能正常访问了。

    要是swap上也没有这个页了，此时还分两种情况：如果这个页是合法分配的且物理内存不满，OS会在物理内存上分配一个新的页来存储这个数据；如果这个页是非法访问的，OS就会认为这是一个严重的错误，可能是程序访问了一个未分配的内存地址或者访问了一个受保护的内存地址，此时报错 \texttt{segmentation fault} ，且OS会强行终止这个进程。这大概率是程序出了bug，访问了一个非法的地址；也有可能是这是个恶意程序（但被OS发现了）。

    如果物理内存和swap都满了，OOM Killer就会自动杀掉一些占用内存过多的进程来释放内存，保证系统的稳定性。

\end{frame}

\begin{frame}
    \frametitle{内存分配器工作原理}

    内存分配器是操作系统中负责管理内存分配和回收的组件，最常见的莫过于 \texttt{malloc} 和 \texttt{free} 了。内存分配器的工作原理可以概括为以下几个步骤：
    \begin{enumerate}
        \item 内存分配器维护一个内存池，通常是从操作系统请求的一大块内存。
        \item 当程序调用 \texttt{malloc} 请求分配内存时，内存分配器会在内存池中寻找一块足够大的空闲内存块来满足请求。
        \item 如果找到合适的空闲块，内存分配器会将其标记为已分配，并返回该块的地址给程序。
        \item 当程序调用 \texttt{free} 释放内存时，内存分配器会将对应的内存块标记为可用，并可能会将相邻的空闲块合并成更大的块，以减少内存碎片。
    \end{enumerate}

\end{frame}

\begin{frame}
    \frametitle{实际上的内存分配}

    但上述分配器是非常朴素的，实际上的内存分配器会有更多的优化和复杂性。

    例如我们在实际情况下 \texttt{malloc} 10GB内存，会发现电脑内存用量并没有增加10GB。这是因为现代操作系统使用了一个叫做“延迟分配”（lazy allocation）的技术，虽然程序员看来确实是获得了10GB的内存，但实际上这10GB是用到才给。只有少数情况下，会预先分配少量的物理页。

    换句话说，这玩意分配的是虚拟内存（不是物理内存）！

\end{frame}

\subsection{怎么压榨CPU}

\begin{frame}
    \frametitle{缓存}

    我们知道，CPU的速度很快，内存的速度也很快，但这两个东西相互通信要过一个0.1毫秒的内存总线。但0.1毫秒对CPU来说是一个非常长的时间，这些时间中CPU难道只能干瞪眼吗？当然不是，CPU有一个叫做缓存（cache）的东西，来弥补CPU和内存之间的速度差距。

    缓存是一块小的存储器，位于CPU内部或与CPU紧密相连，访问速度非常快。通常分三级：约1ns的L1，约4ns的L2，约10ns的L3。这些缓存的容量依次增大。相比之下，内存的访问延迟大约是100ns。

    CPU会把最近用过的东西和将要用的东西放在缓存里，这样就能快速访问了。CPU还会使用一些算法来预测程序接下来可能会访问哪些数据，并提前把它们加载到缓存中，这就是所谓的“预取”（prefetching）。

\end{frame}

\begin{frame}
    \frametitle{缓存行}

    和“页”类似，缓存的基本单位是“行”。每次把数据从内存加载到缓存中，都是成行搬的。
    
    每一个缓存行都有两个标签，一个记录该缓存行对应的内存地址的高位部分，另一个是valid位记录该缓存行是否有效。

    当CPU试图访问一个内存地址时，会先去L1缓存中查找这段内存有没有被缓存。如果有就直接拿来算，如果没有就“缓存不命中”。
\end{frame}

\begin{frame}
    \frametitle{不命中常见工作原理}

    缓存不命中是难免的。读不命中指的是试图读数据时发生的缓存不命中。
    
    当发生读不命中时，CPU就会逐级在更低级的缓存中查找，如果在L2或L3中找到，就把数据搬到L1中；如果在所有级别的缓存中都没有找到，就只能从内存中把数据搬到L1中了。

    如果L1满了，那就不得不选择一个缓存行来替换掉。通常使用LRU或FIFO等算法来选择被替换的缓存行。如果这个缓存行被修改过了（即dirty），还需要把它写回内存中，这样就会增加额外的开销。

    写不命中指的是试图写数据时发生的缓存不命中。同理，会先试图寻找数据；而在写入数据的时候，一般也有一些策略来决定是直接写入内存（write-through）还是先写入缓存（write-back）。如果是write-back，那么当缓存行被替换掉时，才会把修改过的数据写回内存；如果是write-through，那么每次写数据时都会同时更新内存和缓存。

\end{frame}


\begin{frame}
    \frametitle{局部性原理}

    那么问题就变成了，我们要怎么做，才能减少缓存不命中、防止搬运呢？

    只要把经常一起用的数据放在一个连续的缓存行上，就能一口气全带走了。然后，在搬进来下一个数据之前，先把这个数据尽可能用完，这样就能最大程度地利用缓存了。这就是所谓的“局部性原理”。

    局部性原理分两大方面：
    \begin{description}
        \item[时间局部性] 如果一个数据被访问了，那么它很可能在不久的将来再次被访问。
        \item[空间局部性] 如果一个数据被访问了，那么它附近的数据也很可能在不久的将来被访问。 
    \end{description}
    利用局部性原理，我们可以通过一些优化手段来提高程序的性能，最常见的就是数组按行优先遍历（row-major order）而不是按列优先遍历（column-major order），因为前者更符合空间局部性。
\end{frame}

\begin{frame}
    \frametitle{组相联}

    除了利用局部性原理，组相联也是一种常见的缓存组织方式。在这种方式下，缓存被划分为若干个组，每个组包含若干个缓存行。当CPU访问一个内存地址时，MMU会根据地址计算出该地址应该映射到哪个组，然后在该组内查找是否有匹配的缓存行。

\end{frame}

\begin{frame}
    \frametitle{流水线}

    上述缓存机制虽然显著提升了CPU的性能，但是CPU依然有一个瓶颈：指令执行的速度远远跟不上CPU的时钟频率。时钟频率指的是CPU时间的最小单位。

    为了克服这个瓶颈，现代CPU采用了流水线（pipeline）技术，将指令的执行过程分解为多个阶段，使得多个指令可以同时在不同的阶段执行，从而提高CPU的吞吐量。这个流水线和工厂内的流水线很类似，把指令的执行过程分成多个阶段，每个阶段由一个专门的硬件单元来处理。当一个指令在某个阶段执行时，下一条指令就可以进入下一个阶段执行，这样就能实现指令的并行执行了。

    这五个阶段分别是：取指（IF）、译码（ID）、执行（EX）、访存（MEM）和写回（WB）。前一条指令执行到执行阶段，下一个指令就发射。于是每一个时钟周期，CPU都能完成一条指令的执行了，极大地提升了性能。

\end{frame}

\begin{frame}
    \frametitle{冒险及其解决流程}

    但流水线不会一成不变地执行下去，有时候会出现一些情况，导致流水线无法正常工作，这些情况被称为“冒险”（hazard）。

    一般有三类冒险：
    \begin{description}
        \item[数据冒险] 当一个指令需要使用前一条指令的结果时，如果前一条指令还没有完成，那么就会发生数据冒险。处理方法：数据旁路，或插入气泡。避免方法：减少数据依赖（如全局变量等）。
        \item[控制冒险] 当一个分支指令（如if语句）被执行时，CPU需要决定下一条指令的地址，但如果分支条件还没有计算出来，那么就会发生控制冒险。处理方法：分支预测（猜），或插入气泡。避免方法：减少分支指令，保持分支可预测。
        \item[结构冒险] 当两个或更多的指令需要使用同一个硬件资源（如寄存器、ALU等）时，如果这个资源不能同时满足这些指令的需求，那么就会发生结构冒险。处理方法：增加资源，或插入气泡。避免方法：使用小而整齐的循环体，减少资源竞争。
    \end{description}

\end{frame}

\begin{frame}
    \frametitle{更进一步}

    更进一步压榨CPU的手段还有很多：
    \begin{description}
        \item[超标量] 在每个时钟周期内，CPU可以同时发射多条指令到不同的执行单元中执行，从而进一步提高指令级并行性。
        \item[OoO] 指令不必严格按照程序的顺序执行，只要数据够了就可以抢跑。
        \item[超线程] 一个流水线塞两个线程的指令，来提高CPU的利用率。
        \item[big.LITTLE] 在同一颗CPU上塞少量的大核（乱序宽发高功耗）、大量的小核（顺序窄发低功耗），来兼顾性能和功耗。
        \item[向量化] 同时处理多个数据元素的指令，来提高数据级并行性。 
        \item[增加内存带宽] 朴素得有些可笑，但确实是一个极为重要的手段。
    \end{description}

\end{frame}

\subsection{系统调用}

\begin{frame}
    \frametitle{无奈之举}

    刚刚我们说过，现代OS下的用户态软件只能访问自己的虚拟地址空间，不能访问其他进程的内存，也不能直接操作硬件资源（如磁盘、网络等）。但切实的问题是：有的用户态软件他就是要访问这些资源的啊！怎么办呢？

    这就需要一个叫做系统调用（system call）的机制，来让用户态软件能够请求OS提供服务。实际上我们可以把系统调用当成用户态程序试图 \texttt{sudo} 。

    当然，\texttt{sudo} 权限不能白给，你得告诉我你要干嘛，我才能决定要不要给你这个权限。系统调用也是一样的，你得告诉OS你要干嘛（干什么、干多少次、怎么干），OS才能决定要不要给你这个权限。

\end{frame}

\begin{frame}
    \frametitle{系统调用常见处理流程}

    以 \texttt{printf("Hello")} 为例，这个东西实际上是做了一次系统调用 \texttt{write(1, buf, 5)} 。
    
    现在glibc把这玩意塞进寄存器触发syscall指令，然后CPU暂停当前代码的执行，把当前执行状态（寄存器、程序计数器等）保存到内核栈中，然后根据系统调用号（在这里是1，代表write）找到对应的系统调用处理函数，并切换到内核态。
    
    之后，CPU在内核态办事：检查文件描述符1是否可写，发现可以写，就把Hello这5个字节写入到文件描述符1对应的设备（通常是终端）。
    
    写完后，CPU会将结果（成功写入的字节数，在这里是5）放回RAX寄存器，然后切换回用户态。
    
    然后代码继续执行。

\end{frame}

\begin{frame}
    \frametitle{系统调用的代价}

    我们看出来了一个问题：虽然系统调用提供了一个安全的接口来访问OS资源，但它的代价也不小。每次系统调用都需要进行用户态和内核态之间的切换，这个过程涉及到保存和恢复CPU状态、切换内存上下文等操作，通常需要几十到几百个CPU周期。

    因此，在性能敏感的代码中，频繁的系统调用大概率会是一个瓶颈，所以这玩意能不用就不用。一定要用的时候，也要少用。例如使用缓冲IO、批量读写等技术来减少系统调用的次数。

\end{frame}

\section{计算机联网}

\begin{frame}
    \frametitle{一些名词}

    我们刚才所有的内容都是在单机上发生的，但现在的计算机已经不是孤立的了，而是通过网络连接在一起的了。我们来看看一些和网络相关的名词：
    \begin{description}
        \item[带宽] 网络的理论最大传输速率，通常以比特每秒（bps）为单位。
        \item[延迟] 数据从源头到目的地的时间，通常以毫秒（ms）为单位。
        \item[传输速率] 数据的实际传输速率，通常以比特每秒（bps）为单位。
        \item[丢包率] 数据包在传输过程中丢失的比例，通常以百分比（\%）为单位。 
    \end{description}

\end{frame}

\begin{frame}
    \frametitle{另一些名词}

    \begin{description}
        \item[IP地址] 每台连接到网络的计算机都有一个唯一的IP地址，用来标识它在网络中的位置。可以把这玩意当成门牌号。分IPv4和IPv6两种，前者是32位的地址，后者是128位的地址。
        \item[MAC地址] 每个网络接口都有一个唯一的MAC地址，用来标识它在局域网中的位置。
        \item[端口号] 每个网络服务都有一个端口号，用来标识它在计算机上的位置。
        \item[局域网（LAN）] 连接在一起的计算机组成的网络，通常覆盖一个小范围，如家庭、办公室等。
        \item[广域网（WAN）] 连接在一起的计算机组成的网络，通常覆盖一个大范围，如互联网。互联网是最大的WAN。 
    \end{description}

\end{frame}

\begin{frame}
    \frametitle{常见网络设备}

    \begin{description}
        \item[调制解调器] 是负责将数字信号转换为模拟信号（调制）和将模拟信号转换为数字信号（解调）的设备，俗称“猫”，通常用于连接到互联网。传统的电话线猫已经被彻底淘汰，现在的猫一般是光纤猫或电缆猫。
        \item[路由器] 连接不同网络的设备，负责转发数据包到正确的目的地。
        \item[交换机] 连接同一网络中的设备，负责转发数据包到正确的设备。
    \end{description}

\end{frame}

\begin{frame}
    \frametitle{常见网络协议}

    网络协议是指在网络通信中使用的一套规则和标准，用于定义数据的格式、传输方式、错误处理等。常见的网络协议有：
    \begin{description}
        \item[HTTP] 超文本传输协议，用于在Web上进行数据传输。
        \item[FTP] 文件传输协议，用于在网络上进行文件传输。
        \item[SMTP] 简单邮件传输协议，用于在网络上发送电子邮件。
        \item[TCP] 传输控制协议，提供可靠的、面向连接的通信服务。
        \item[UDP] 用户数据报协议，提供不可靠的、无连接的通信服务。
    \end{description}

\end{frame}

\begin{frame}
    \frametitle{内网和外网}

    内网是相对外网而言的。内网是指在局域网内部使用的IP地址，通常是私有IP地址（如192.168.x.x、10.x.x.x等），这些地址仅仅在内网上有效，无法直接访问互联网。外网是指在互联网上使用的IP地址，通常是公共IP地址，这些地址可以被全球访问。

    网关是连接内网和外网的设备，负责转发数据包到正确的目的地。当内网中的设备要访问某个设备的时候，数据包会先扔给网关；网关检查数据包的目的地址是在直连网段还是外部网络，并发送到正确的地方。

\end{frame}

\begin{frame}
    \frametitle{计算机启动时的网卡上线}

    我们知道，计算机启动时，POST、硬件初始化、操作系统启动流程，内核唤醒驱动。轮到网卡的时候，驱动会给网卡复位、写EEPROM里的MAC地址到网卡寄存器里。于是网卡启动。

    网卡启动后，开始向网线里发广播包，并监听对面有没有回应。对面一般是交换机或路由器，如果一切正常会发个FLP回应。网卡收到回应之后，知道对面有人，接下来就可以开始协商速率、双工模式等了。协商成功后，网卡就在操作系统层面上线了，可以开始正常工作了。

\end{frame}

\begin{frame}
    \frametitle{DORA流程}

    \begin{enumerate}
        \item 网卡上线之后，协议栈还是空的，除了MAC地址啥都没有，无法上网。网卡构造一个以太网广播包，里面装一个DHCP Discover包，表示“我上线了，给我发个IP”，走UDP发走。路由器会检查这个包的目的地址，发现是Discover，于是广播改单播经DHCP中继扔DHCP服务器（现代路由器一般都内置了DHCP服务器）。
        \item DHCP服务器收到包后，查IP池子，找个空闲的IP地址，发个DHCP Offer包回应。该包包含IP地址、地址租期、DNS、网关等信息。
        \item 网卡收到回应后，选一个DHCP服务器，发个Request包表示“我选你了”。
        \item DHCP服务器收到后，发个ACK包表示“OK”，至此网卡就拿到了一个IP地址了。
    \end{enumerate}
    
    整个流程叫做DORA流程（Discover-Offer-Request-Ack）。

\end{frame}

\begin{frame}
    \frametitle{ARP流程}

    现在网卡有了IP地址了，但就像寄信一样，有发件人也要有收件人，这些收件地址都是MAC地址。于是协议栈查路由，发现上述DHCP Offer提供的网关在直连网段，于是构造一个ARP Request包，问：“谁有这个IP地址，给我发个MAC地址来”，走以太网广播发走。

    网关收到该包，回应一个ARP Reply包，告诉网卡这个IP地址对应的MAC地址是什么东西，并缓存60秒。至此，网卡就可以正常上网了。

    ARP还可以防止IP地址冲突。如果多台主机的IP相同，会回复包，提示用户修改IP地址。
\end{frame}

\begin{frame}
    \frametitle{路由流程}

    如果需要访问外网，那流程就长得多了。
    \begin{enumerate}
        \item 协议栈查路由表，发现目的地址不在直连网段，于是走默认路由，构造一个以太网帧，目的地址为网关，源为自己。
        \item 网关（即交换机或路由器）收到包后，检查目的地址，发现不认识这个IP，于是剥掉二层头、查三层转发表，找到下一跳的IP，TTL减一，重新计算校验和，构造一个新的以太网帧，目的地址为下一跳的MAC地址，源被NAT为自己的MAC地址。NAT是一种能将私有的IP地址转换为公共IP地址的技术，通常用于内网访问外网。
        \item 下一跳收到包后，重复上述过程（但不会NAT了），直到包到达目的地。
    \end{enumerate}
    如果在包到达目的地之前就发生了TTL变为0的情况，路由器会发一个ICMP Time Exceeded包回源地址，告诉它这个包被丢弃了。这样能防止包在网络中无限循环。这也正是 \texttt{traceroute} 的原理。

\end{frame}

\begin{frame}
    \frametitle{DNS流程}

    现在万事俱备，理论上可以随意地访问互联网了；但背诵一大堆IP地址也不是一个好办法。为了方便用户，人们发明了DNS系统。DNS本质也是一个从字符串到IP的映射，把易记的域名 \url{www.google.com} 映射到一个IP地址 \texttt{8.8.8.8} 上。

    浏览器输入URL后，计算机先查本地 \texttt{/etc/hosts} 文件和本地DNS缓存，看看认不认识这个域名。如果不认识，就会向DNS服务器发起一个DNS查询请求。DNS服务器收到请求后，会根据域名的层次结构进行递归查询，最终返回对应的IP地址给计算机，过程类似上述路由流程。

    DNS是非常容易受攻击的：DNS劫持、DNS污染等攻击手段可以篡改DNS的查询结果，甚至可以劫持NXDOMAIN响应来实现域名劫持，从而实现钓鱼网站、随地插入广告等恶意行为。于是也就有了类似DoH等加密DNS协议，防止中间人篡改DNS查询结果。

\end{frame}

\begin{frame}
    \frametitle{TCP三次握手流程}

    UDP直接发送数据包虽然快，但不可靠：数据包可能丢失、重复或乱序。TCP提供了可靠的、面向连接的通信服务，确保数据包按顺序到达，并且没有丢失，其实现原理是“三次握手”流程：
    \begin{enumerate}
        \item 第一次握手：客户端发送一个SYN包（同步包）给服务器，表示请求建立连接。
        \item 第二次握手：服务器收到SYN包后，回复一个SYN-ACK包（同步-确认包），表示同意建立连接，并且确认收到客户端的SYN包。
        \item 第三次握手：客户端收到SYN-ACK包后，回复一个ACK包（确认包），表示连接建立成功。
        \item 连接完成后，双方各自维护自己的源IP、源端口、目的IP、目的端口四元组，滑动窗口、序列号、拥塞窗口等状态信息来保证数据传输的可靠性和效率。
    \end{enumerate}
    TCP虽然慢但可靠，适用于需要保证数据完整性的应用，如Web浏览、文件传输等；UDP虽然快但不可靠，适用于对实时性要求较高但对数据完整性要求不高的应用，如视频直播、在线游戏等。

\end{frame}

\begin{frame}
    \frametitle{HTTP流程}

    有了TCP，我们就可以可靠地和服务器相互通信了。但是数据包是二进制的，浏览器看不懂。为了解决这个问题，引入HTTP技术把数据包格式化成文本，便于浏览器等客户端理解和处理。

    浏览器访问某IP时，首先会建立一个TCP连接（如果之前没有的话），然后发送一个HTTP请求包给服务器。服务器收到请求后，处理请求并生成一个HTTP响应包，包含状态码、响应头和响应体等信息，然后发送回客户端。客户端收到响应后，根据状态码和响应头来处理响应体，把经典三件套（HTML、CSS、JavaScript）渲染成网页展示给用户。

    HTTP是无状态的协议，每一个请求都是独立的，服务器不会记住之前的请求。为了实现会话管理，HTTP引入了cookie机制，允许服务器在响应中设置cookie，客户端在后续请求中携带这些cookie，从而实现状态的维护。

\end{frame}

\begin{frame}
    \frametitle{TLS流程}

    HTTP虽然方便，但是是明文的，容易被窃听和篡改。为了保护数据的安全性，引入了TLS（传输层安全协议）来加密HTTP通信，形成HTTPS。

    当客户端访问一个HTTPS URL时，首先会建立一个TCP连接，然后进行TLS握手流程来协商加密算法和密钥。TLS握手流程包括以下步骤：
    \begin{enumerate}
        \item 客户端发送一个ClientHello消息，包含支持的TLS版本、加密算法列表、随机数等信息。
        \item 服务器收到ClientHello后，回复一个ServerHello消息，选择TLS版本和加密算法，并发送服务器证书（包含公钥）给客户端。
        \item 客户端验证服务器证书的合法性，生成一个随机数，并使用服务器的公钥加密这个随机数，发送给服务器。
        \item 服务器收到加密的随机数后，使用自己的私钥解密得到随机数，然后双方使用这个随机数和之前协商的加密算法来生成对称密钥，后续的通信都使用这个对称密钥进行加密。
    \end{enumerate}
    TLS握手完成后，客户端和服务器就可以使用协商好的加密算法和密钥来加密HTTP请求和响应了，从而保护数据的安全性。当然TLS也很慢，所以也一直是优化的重点。

\end{frame}

\begin{frame}
    \frametitle{WiFi原理}

    上述内容实际上都是在有线的基础上讲的，无线网络也大差不差，只不过多了一个无线接入点（AP）来充当交换机的角色，负责转发无线设备之间的数据包。

    大致流程是：
    \begin{enumerate}
        \item STA扫描AP，找到可用的AP列表。
        \item STA选择一个AP，发送一个Association Request包，表示要连接这个AP。
        \item AP收到请求后，回复一个Association Response包，表示同意连接。
        \item 连接完成后，STA和AP之间就可以开始正常通信了。
    \end{enumerate}
    之后的流程和有线网络几乎一致，只不过在物理层上使用了无线信号来传输数据包。WiFi还引入了很多优化和安全机制，如WPA2、WPA3等来保护无线通信的安全性。

\end{frame}

\begin{frame}
    \frametitle{CDN原理}

    相对现代数据量而言，网络还是太慢了。为了提升用户访问的速度和体验，CDN（内容分发网络）应运而生了。这东西完全可以理解为“网络上的高速缓存”。

    例如在北京使用淘宝，图片在上海走杭州到北京，延迟会非常高。CDN则会把图片缓存到北京或附近的边缘服务器（天津、河北等）上，用户走DNS解析就会解析到CDN上的服务器上，这样就能大幅降低访问延迟了。

\end{frame}

\begin{frame}
    \frametitle{启示：现代网络明显的分层设计}

    由上述讨论，我们可以看出现代网络明显分为五个层次：
    \begin{description}
        \item[物理层] 指实际的物理介质，如光纤、网线、无线信号等。
        \item[数据链路层] 负责在物理层上建立和维护数据传输的链路，如以太网协议、WiFi协议等。
        \item[网络层] 负责数据包的路由和转发，如IP协议等。
        \item[传输层] 负责数据的传输，如TCP协议、UDP协议等。
        \item[应用层] 负责具体的应用协议，如HTTP协议、DNS协议等。 
    \end{description}
    我们看到的都是应用层的东西，但这些东西并不是空中楼阁，而是建立在下面的层次之上的。每一层都有自己的协议和机制，来实现特定的功能，并且通过接口与其他层进行交互。

    

\end{frame}

% \begin{frame}
%     \frametitle{代理和VPN原理}

%     内网像一座有围墙的城池，出来很容易，进去很困难。但需求是存在的：有时候，我们在校外打比赛或休假，想要访问学校内网的资源，怎么办？这就需要代理和VPN的帮助了。

%     代理是一种中间人服务，即请一位“第三者服务器”来替我们访问内网资源。我们把请求发给这个服务器，它再帮我们访问内网资源，然后把结果返回给我们。常见的代理协议有很多，如HTTP代理、SOCKS代理等。

%     VPN则不需要第三者服务器，而是直接在用户设备和内网之间建立一个加密的隧道，使得用户设备看起来就像直接连接在内网一样。这样用户就可以直接访问内网资源了。常见的VPN协议有PPTP、L2TP、OpenVPN等。

%     举例：北京大学VPN服务。

% \end{frame}

\end{document}